\documentclass[12pt]{article}
\usepackage[T1]{fontenc}
\usepackage[utf8]{inputenc}
\usepackage{lmodern}
\usepackage{textcomp}
\usepackage{enumitem}
\usepackage{geometry}
\usepackage{graphicx}
\usepackage{amsmath, amssymb}
\geometry{margin=1in}
\begin{document}
\begin{center}
Astronomy - Graduate Level Assessment 3 \
Total Marks: 40 \
Answer all questions.
\end{center}

\section*{Multiple Choice (10 marks)}
\begin{enumerate}[label=\arabic*.]
\item The Large Magellanic Cloud is ...
\begin{enumerate}[label=(\alph*)]
    \item a dwarf galaxy orbiting the Milky Way.
    \item the closest planetary nebula to the Earth.
    \item a bright star cluster discovered by Magellan.
    \item the outer arm of the Milky Way named after Magellan.
\end{enumerate}
\item Which factor s most important in determining the history of volanism and tectonism on a planet?
\begin{enumerate}[label=(\alph*)]
    \item size of the planet
    \item presence of an atmosphere
    \item distance from the sun
    \item rotation period
\end{enumerate}
\item Which of the following most likely explains why Venus does not have a strong magnetic field?
\begin{enumerate}[label=(\alph*)]
    \item Its rotation is too slow.
    \item It has too thick an atmosphere.
    \item It is too large.
    \item It does not have a metallic core.
\end{enumerate}
\item The Pleiades is an open star cluster that plays a role in many ancient stories and is well-known for containing ... bright stars.
\begin{enumerate}[label=(\alph*)]
    \item 5
    \item 7
    \item 9
    \item 12
\end{enumerate}
\item Why is the sky blue?
\begin{enumerate}[label=(\alph*)]
    \item Because the molecules that compose the Earth's atmosphere have a blue-ish color.
    \item Because the sky reflects the color of the Earth's oceans.
    \item Because the atmosphere preferentially scatters short wavelengths.
    \item Because the Earth's atmosphere preferentially absorbs all other colors.
\end{enumerate}
\end{enumerate}

\section*{True / False (10 marks)}
\textit{Answer True or False}
\begin{enumerate}[label=\arabic*.]
\setcounter{enumi}{5}
\item True or False: Hydrogen is the most abundant element found in the atmospheres of the Jovian planets.
\item True or False: Mars experiences more extreme seasons due to its more eccentric orbit compared to Earth's nearly circular orbit.
\item True or False: The atmosphere on Mars is too thin to effectively trap a significant amount of heat, limiting the greenhouse effect.
\item True or False: Moons contribute to the stability of particles within rings, create gravitational effects at ring edges, and form gaps between rings.
\item True or False: When the angular separation of two stars is smaller than the angular resolution of the eye, they will be distinctly visible as two separate stars.
\end{enumerate}

\section*{Long Form Response (20 marks)}
\textit{Answer each question with a clear and complete explanation.}
\begin{enumerate}[label=\arabic*.]
\setcounter{enumi}{10}
\item Discuss the phenomena associated with orbital resonance in astronomy, and analyze why the breaking of small Jovian moons to form ring materials is not a consequence of this process.
\item Discuss the concept of a black body in astronomy, including its definition, significance in radiation calculations, and implications for understanding stellar properties.
\end{enumerate}

\vfill
\noindent\textit{End of Paper}
\end{document}